% Created (UTC): 2026-06-21T04:03:14Z
% Repository HEAD: abd2e7116aa9b8eeaf0206507ec9181346bf65b3
% Companion to: src/PolyLog/docs/article/gaussian-eisenstein-double-polylogs.tex
\documentclass[11pt]{article}
\usepackage[margin=1in]{geometry}
\usepackage{amsmath,amssymb,amsthm,mathtools}
\usepackage{booktabs}
\usepackage{microtype}
\emergencystretch=3em
\hbadness=4000
\usepackage[colorlinks=true,linkcolor=blue,citecolor=blue,urlcolor=blue]{hyperref}

\DeclareMathOperator{\Li}{Li}
\DeclareMathOperator{\Cl}{Cl}
\DeclareMathOperator{\Real}{Re}
\DeclareMathOperator{\Imag}{Im}
\DeclareMathOperator{\lcm}{lcm}
\DeclareMathOperator{\ord}{ord}
\newcommand{\iu}{\mathrm{i}}
\newcommand{\rr}{\rho}
\newcommand{\code}[1]{\texttt{\small #1}}

\theoremstyle{plain}
\newtheorem{theorem}{Theorem}
\newtheorem{proposition}[theorem]{Proposition}
\theoremstyle{remark}
\newtheorem{remark}[theorem]{Remark}

\title{Completing the level-3 and level-4 frontier:\\
the $w=1$ mixed cube-root double polylogarithms}
\author{Vladimir Reshetnikov}
\date{2026-06-21}

\begin{document}
\maketitle

\begin{center}
\small
Research report. Created (UTC) 2026-06-21T04:03:14Z; repository \texttt{HEAD}
\texttt{abd2e7116aa9b8eeaf0206507ec9181346bf65b3}.\\
Companion to \emph{Multiple polylogarithms at Gaussian and Eisenstein arguments}
(\texttt{gaussian-eisenstein-double-polylogs}).
\end{center}

\begin{abstract}
The article \emph{Multiple polylogarithms at Gaussian and Eisenstein arguments} left one piece
unfinished: the \emph{mixed} cube-root points $\Li_{a,b}(\rr^2,\rr)$ (and the Gaussian
$\Li_{a,b}(\iu,-\iu)$), where $w=z_1z_2=1$ degenerates the accelerated evaluator. This report builds a
robust $w=1$ evaluator---a smooth-correction identity accelerated by Wynn's $\rho$-algorithm---generates
the previously inaccessible values to $500$ digits, reduces them to explicit closed forms through
weight~$4$, and completes the full level-3 and level-4 depth-2 double-shuffle dimension. The headline:
in \emph{both} towers the mixed $w=1$ point adds, beyond the article's survivors, exactly \emph{two} new
depth-2 constants---a real $\Real\Li_{4,1}$ direction at weight~$5$ and an imaginary $\Imag\Li_{5,1}$
direction at weight~$6$---in identical positions.
\end{abstract}

\section{The one remaining piece}\label{sec:intro}

The companion article closed its level-3 (Eisenstein) and level-4 (Gaussian) double-shuffle dimension
computations with one explicitly flagged gap. In the Eisenstein run the \emph{mixed} cube-root points
\[
  \Li_{a,b}(\rr^2,\rr),\qquad \Li_{a,b}(\rr,\rr^2),\qquad \rr=e^{2\pi\iu/3},
\]
were never folded in, because there $w=z_1z_2=\rr^3=1$, and (quoting the article)
\begin{quote}\itshape
both the Lerch-plus-Wynn evaluator and a direct cumulative-sum evaluator stall there at ${\sim}8$
digits---the period-three, non-oscillating tail defeats Wynn's $\varepsilon$-algorithm \dots\ Folding
the mixed cube-root points in (for the full level-3 frontier) requires a period-three-aware accelerated
evaluator, or Wolfram; it is the one remaining piece.
\end{quote}
The Gaussian run skipped its own $w=1$ point $(\iu,-\iu)$ for the same reason. This report builds the
missing evaluator, generates the previously inaccessible values to several hundred digits, and uses
them to (i)~harvest closed forms for the mixed-point values and (ii)~complete the full level-3 and
level-4 depth-2 double-shuffle dimension.

Throughout, the convention is
\begin{equation}\label{eq:conv}
  \Li_{a,b}(z_1,z_2)=\sum_{n_1>n_2\ge1}\frac{z_1^{n_1}z_2^{n_2}}{n_1^{a}\,n_2^{b}},
\end{equation}
which is \code{MultiplePolyLog[\{a,b\},\{z1,z2\}]} in Wolfram~15 (exponents first).

\section{Why the standard evaluator degenerates at \texorpdfstring{$w=1$}{w=1}}\label{sec:why}

With the convention \eqref{eq:conv}, the accelerated evaluator of the article writes the inner tail as
a Lerch transcendent and sums the \emph{outer} series with an acceleration factor $w^n=(z_1z_2)^n$. At
the mixed points $w=1$, that factor is identically $1$: the oscillation Wynn's $\varepsilon$-algorithm
relies on disappears, leaving a bare period-three power-law tail that the algorithm cannot resolve past
a handful of digits. Summing over the \emph{larger} index instead keeps an oscillating factor
$z_1^{n_1}$, but for $z_1=\rr^2$ this is a \emph{period-three} oscillation, and a plain
$\varepsilon$-algorithm (tuned for geometric tails) again stalls.

\section{A robust \texorpdfstring{$w=1$}{w=1} evaluator}\label{sec:eval}

Two ingredients fix it.

\paragraph{(a) A smooth correction identity.} Writing $l=n_1$ (the larger index) and splitting the
inner partial sum $P(l)=\sum_{n<l}z_2^n/n^b=\Li_b(z_2)-R(l)$ with tail
$R(l)=\sum_{n\ge l}z_2^n/n^b=z_2^{\,l}\,\Phi(z_2,b,l)$ (Lerch $\Phi$), one gets the exact identity
\begin{equation}\label{eq:smooth}
  \boxed{\;\Li_{a,b}(z_1,z_2)=\Li_a(z_1)\,\Li_b(z_2)\;-\;\sum_{l\ge1}\frac{w^{\,l}}{l^{a}}\,\Phi(z_2,b,l)\;}
\end{equation}
At $w=1$ the correction summand $\Phi(z_2,b,l)/l^a$ is \emph{smooth}---every oscillating factor
$z_1^l z_2^l=w^l$ collapses to $1$---with the clean asymptotic
$\Phi(z_2,b,l)\sim\sum_{j\ge0}\binom{-b}{j}c_j(z_2)\,l^{-b-j}$, where $c_0=\tfrac1{1-z_2}$ and
$c_{j\ge1}=\Li_{-j}(z_2)$ (rational in $z_2$). The subtle point---the source of a factor-of-one bug
that cost five digits in a first draft---is that $c_0=\sum_{k\ge0}z_2^k=\tfrac1{1-z_2}$ is \emph{not}
$\Li_0(z_2)=\tfrac{z_2}{1-z_2}$; they differ by exactly $1$. This gives an Euler--Maclaurin tail in
Hurwitz zetas, $\sum_{l>M}\Phi/l^a=\sum_j\binom{-b}{j}c_j\,\zeta(a+b+j,M+1)$, a fast and accurate
cross-check at low to medium precision.

\paragraph{(b) Period-aware $\rho$-acceleration for production.} For high precision the production
evaluator sums the outer cumulative series directly, samples the partial sums at the common period
$N=\lcm(\ord z_1,\ord z_2)$ (so every oscillating factor realigns to $1$ at the sample points), and
accelerates the resulting smooth, power-law-convergent sequence with \textbf{Wynn's $\rho$-algorithm}
---the algebraic-convergence counterpart of the $\varepsilon$-algorithm. The $\rho$-algorithm yields
roughly one correct digit per sample point; the evaluator inflates its internal precision to absorb the
conditioning loss and returns the requested number of correct digits
(\code{mpl\_w1.py}).

\section{Validation}\label{sec:valid}

The evaluator is checked by three independent, Wolfram-free routes plus one Wolfram cross-check.
\begin{enumerate}
\item \textbf{Exact symmetric closed form} (single polylogs only, an entirely separate code path):
  \begin{equation}\label{eq:sym}
    \Li_{a,b}(\rr^2,\rr)+\Li_{b,a}(\rr,\rr^2)=\Li_a(\rr^2)\Li_b(\rr)-\zeta(a+b),
  \end{equation}
  and the Gaussian analogue $\Li_{a,b}(\iu,-\iu)+\Li_{b,a}(-\iu,\iu)=\Li_a(\iu)\Li_b(-\iu)-\zeta(a+b)$.
  (The collision term is $\Li_{a+b}(z_1z_2)=\Li_{a+b}(1)=\zeta(a+b)$.)
\item \textbf{Conjugation} $\Li_{a,b}(\rr,\rr^2)=\overline{\Li_{a,b}(\rr^2,\rr)}$.
\item \textbf{Two independent summation methods}---the $\rho$-accelerated cumulative sum and the
  Euler--Maclaurin smooth-correction series---agree to full working precision.
\item \textbf{Wolfram~15} \texttt{MultiplePolyLog} at moderate precision.
\end{enumerate}
All four agree; the symmetric closed form \eqref{eq:sym} holds to the full generated precision
(${\sim}10^{-500}$) for every $(a,b)$.

\section{Closed forms for the Eisenstein mixed-point values}\label{sec:eisforms}

With $500$-digit values in hand, an integer-relation (PSLQ/LLL) search against the level-3 basis
---monomials in the single $L$-values $\{\pi,\log 3,\zeta(3),\zeta(5),\Cl_2(\tfrac{2\pi}3),
\Cl_4(\tfrac{2\pi}3),\Cl_6(\tfrac{2\pi}3)\}$ times the $(\rr,1)$ generators
$g_{a,b}=\Imag\Li_{a,b}(\rr,1)$, $h_{a,b}=\Real\Li_{a,b}(\rr,1)$, and the lone $(\rr,\rr)$
double-shuffle survivor $\Imag\Li_{3,2}(\rr,\rr)$---\textbf{reduces every mixed value through weight
$4$, and every weight-5 imaginary part}, to explicit closed forms (each canary-clean, height $\le 10$,
re-verified to $10^{-425}$ against the \emph{complete} weight-graded basis). The two simplest:
\begin{equation}\label{eq:eis-w2}
  \Imag\Li_{1,1}(\rr^2,\rr)=-\bigl(\Cl_2(\tfrac{2\pi}3)+\Imag\Li_{1,1}(\rr,1)\bigr),\qquad
  \Real\Li_{1,1}(\rr^2,\rr)=5\,\Real\Li_{1,1}(\rr,1)-\tfrac12\log^2 3.
\end{equation}
(The full list of ${\sim}16$ reductions is in \code{harvest\_eis\_mixed.py}.) The weight-5 \emph{real}
parts and all of weight $6$ do \emph{not} reduce to this basis at height $\lesssim 10^{4}$.

\section{The completed level-3 double-shuffle dimension}\label{sec:eisdim}

Two computations pin the dimension of the full level-3 depth-2 space---all four colored points
$(\rr,\rr),(\rr^2,\rr^2),(\rr^2,\rr),(\rr,\rr^2)$---modulo single polylogarithms and the $(\rr,1)$
generators.
\begin{itemize}
\item The \textbf{exact-rank convergent double shuffle} (stuffle and shuffle of products
  $\Li_p(s)\Li_q(t)$, $s,t\in\{\rr,\rr^2\}$, together with the convergent zeta-products
  $\Li_p(s)\zeta(q)$ whose shuffle reaches the mixed points) gives an \emph{upper bound} of $\mathbf 7$
  through weight $6$ (the divergent $\Li_1$ words, needing shuffle regularization, are omitted).
\item \textbf{Combining} that exact rank with the PSLQ closed forms of Section~\ref{sec:eisforms}
  collapses it:
\end{itemize}
\begin{center}
\begin{tabular}{ccl}
\toprule
weight & irreducible dimension & survivors\\
\midrule
$2$ & $0$ & ---\\
$3$ & $0$ & ---\\
$4$ & $0$ & ---\\
$5$ & $2$ & $\Imag\Li_{3,2}(\rr,\rr)$ \emph{(article)},\quad $\Real\Li_{4,1}(\rr^2,\rr)$ \textbf{(new)}\\
$6$ & $1$ & $\Imag\Li_{5,1}(\rr^2,\rr)$ \textbf{(new)}\\
\bottomrule
\end{tabular}
\end{center}
So the full level-3 depth-2 space, mod single polylogs and the $(\rr,1)$ generators, has dimension
$\mathbf 3$ through weight $6$. Everything through weight $4$ collapses entirely into the
$(\rr,1)$-generator world; the mixed cube-root points contribute \textbf{exactly two new constants}
beyond the article's single $(\rr,\rr)$ survivor---a \emph{real} direction $\Real\Li_{4,1}(\rr^2,\rr)$
at weight $5$ and an \emph{imaginary} direction $\Imag\Li_{5,1}(\rr^2,\rr)$ at weight $6$.

\begin{remark}[the article's own lesson]
Non-reduction by integer-relation search is strong evidence, not proof: a genuine reduction of large
height is invisible to PSLQ, and the convergent double shuffle omits the regularized $\Li_1$ relations.
So the two new survivors are best read as \emph{candidate} new level-3 depth-2 constants; a full
regularized double shuffle, or a Bloch-group / Zagier-conjecture argument, would settle their
independence definitively. As an additional check, an \emph{enriched}-basis probe
(\code{probe\_eis\_candidates.py})---adding $\Li_n$ at candidate level-3 CM half-points
$\tfrac{1\pm\rr}2,\ \tfrac1{1-\rr},\ \tfrac{\rr}{\rr-1},\ 1-\rr$ (the level-3 analogue of the Gaussian
$\tfrac{1+\iu}2$ that resolved the article's Family-2)---fails to crack either candidate:
$\Real\Li_{4,1}(\rr^2,\rr)$ resists a $105$-atom basis (spurious scale ${\sim}10^{4}$).
\end{remark}

\section{The Gaussian \texorpdfstring{$(\iu,-\iu)$}{(i,-i)} point and an exact parallel}\label{sec:gauss}
\sloppy
The same evaluator handles the level-4 Gaussian point $(\iu,-\iu)$ that the article's level-4 run also
skipped ($w=\iu\cdot(-\iu)=1$). The picture is identical in shape.
\begin{itemize}
\item \textbf{Closed forms.} PSLQ against the level-4 basis (single $L$-values
  $\pi,\log 2,\beta(2),\beta(4),\beta(6),\zeta(3),\zeta(5)$, the half-point values
  $\Real\Li_n(\tfrac{1+\iu}2),\Imag\Li_n(\tfrac{1+\iu}2)$, and the $(\iu,1)$ generators) reduces every
  $(\iu,-\iu)$ value through weight $3$ and every weight-4 imaginary part. The weight-2 case is
  especially clean:
  \begin{equation}\label{eq:gauss-w2}
    \Li_{1,1}(\iu,-\iu)=-\Li_2\!\left(\tfrac{1-\iu}2\right)\qquad(\text{verified to } 10^{-50}).
  \end{equation}
\item \textbf{Dimension.} The corrected exact-rank double shuffle (now including $(\iu,-\iu)$) gives an
  upper bound of $12$; combined with the closed forms it collapses to the article's five $(\iu,-1)$
  constants plus \textbf{the same two new mixed directions}---$\Real\Li_{4,1}(\iu,-\iu)$ at weight $5$
  and $\Imag\Li_{5,1}(\iu,-\iu)$ at weight $6$. (One weight-3 residual, $\Imag\Li_{2,1}(\iu,-1)$, is a
  value the article already reduces to the half-point $\lambda/\mu$ basis, outside the convergent
  double shuffle's reach.)
\end{itemize}

\begin{theorem}[the parallel]\label{thm:parallel}
In \emph{both} towers the $w=1$ mixed point contributes, beyond the published survivors, exactly two
new depth-2 constants, and \emph{in the same two positions}:
\[
  \boxed{\ \Real\Li_{4,1}(\mathrm{mixed})\ \text{at weight }5,\qquad
            \Imag\Li_{5,1}(\mathrm{mixed})\ \text{at weight }6.\ }
\]
\end{theorem}
That the level-3 and level-4 mixed points add new content at identical weights and index patterns
---a \emph{real} $\Li_{4,1}$ direction and an \emph{imaginary} $\Li_{5,1}$ direction---is the report's
most suggestive finding, and a natural target for the regularized double shuffle to explain.

\section{Significance}\label{sec:sig}

\begin{itemize}
\item \textbf{The ``one remaining piece'' is closed.} The mixed cube-root points the article could not
  evaluate are now computable to arbitrary precision; the level-3 frontier is complete through weight
  $6$, and the same evaluator handles the Gaussian $(\iu,-\iu)$ point the level-4 run skipped.
\item \textbf{A reusable tool.} The smooth-correction $+$ period-aware $\rho$-acceleration evaluator
  handles \emph{every} root-of-unity double polylogarithm, including the degenerate $w=1$ diagonal where
  Wolfram's own \texttt{MultiplePolyLog} is slow and erratic (fast at low weight, stalling at higher
  weight). It is the natural companion to the article's accelerated evaluator for the oscillating case.
\item \textbf{A clean structural picture.} Through weight $4$ the mixed cube-root doubles carry
  \emph{no} new content; the first genuinely-new mixed constant appears at weight $5$. This mirrors
  ---and sharpens---the article's finding that the colored towers are far smaller than a
  precision-limited search suggests.
\end{itemize}

\section*{Reproducibility}
\sloppy
All scripts are under \code{src/Hypergeometric/tools/scratch/depth3-gaussian/}: \code{mpl\_w1.py}
(evaluator), \code{gen\_eis\_mixed.py} / \code{gen\_gauss\_mixed.py} (value stores \code{vals\_eis\_mixed.py}
/ \code{vals\_gauss\_mixed.py}), \code{harvest\_eis\_mixed.py} / \code{harvest\_gauss\_mixed.py} (closed
forms), \code{shuffle\_ds\_eis\_full.py} / \code{shuffle\_ds\_gauss\_full.py} (exact-rank double shuffle),
and \code{combined\_dim\_eis.py} / \code{combined\_dim\_gauss.py} (the combined dimension). Values are
generated to $500$ digits; every displayed identity is verified to at least $10^{-425}$, and the
symmetric closed form \eqref{eq:sym} (an independent single-polylog identity) holds to ${\sim}10^{-500}$
for every value. The verified closed forms and the dimension results are collected in the
\texttt{Get[]}-able store \code{src/PolyLog/identities/discovered-mixed-cuberoot-doubles.wl}.

\end{document}
